\documentclass[11pt]{article}
\usepackage{inote}
\inotesetup{mode=practice, answers=true, qr=false,
  subject={পদার্থবিজ্ঞান ১ম পত্র}, chapter={অধ্যায় ০২ঃ ভেক্টর}, video=XXXXXXXXXXX}
\begin{document}
\begin{multicols*}{2}
\sheettitle{অধ্যায়-২}{ভেক্টর}
\topicbar{Topic-1: ভেক্টরের লব্ধি}
\problem[লেকচার ০১ · মৌলিক]{৬ একক ও ৮ একক মানের দুটি ভেক্টর একটি বিন্দুতে পরস্পর $60^\circ$ কোণে ক্রিয়া করছে। লব্ধির মান ও দিক নির্ণয় করো।}{$12.17$ একক; ৬ একক ভেক্টরের সাথে $34.71^\circ$}
\problem[লেকচার ০১ · মৌলিক]{$3\,\text{N}$ ও $4\,\text{N}$ মানের দুটি বল পরস্পর লম্বভাবে ক্রিয়া করলে লব্ধির মান ও দিক কত?}{$5\,\text{N}$; $3\,\text{N}$ বলের সাথে $53.13^\circ$}
\problem[লেকচার ০১ · কনসেপ্ট]{সমান মানের দুটি বল $F$ এর লব্ধির মানও $F$ হলে বল দুটির মধ্যবর্তী কোণ কত?}{$120^\circ$}
\problem[লেকচার ০১ · অনুরূপ]{$10\,\text{N}$ ও $20\,\text{N}$ মানের দুটি বল কত কোণে ক্রিয়া করলে লব্ধি $10\,\text{N}$ বলের সাথে সমকোণে থাকবে?}{$120^\circ$}
\topicbar{Topic-2: উপাংশ}
\problem[লেকচার ০১ · মৌলিক]{অনুভূমিকের সাথে $30^\circ$ কোণে ক্রিয়ারত $20\,\text{N}$ বলের অনুভূমিক ও উল্লম্ব উপাংশ নির্ণয় করো।}{$F_x = 17.32\,\text{N}$, $F_y = 10\,\text{N}$}
\topicbar{Topic-3: আপেক্ষিক বেগ}
\problem[লেকচার ০১ · নদী-নৌকা]{$1\,\text{km}$ প্রশস্ত নদীতে স্রোতের বেগ $3\,\text{km\,h}^{-1}$। স্থির পানিতে $5\,\text{km\,h}^{-1}$ বেগের নৌকা সবচেয়ে কম দূরত্বে নদী পার হতে চাইলে কোন দিকে চালাতে হবে এবং কত সময় লাগবে?}{স্রোতের সাথে $126.87^\circ$ কোণে; $15$ মিনিট}
\problem[লেকচার ০১ · বৃষ্টি-ছাতা]{বৃষ্টি খাড়া নিচের দিকে $6\,\text{m\,s}^{-1}$ বেগে পড়ছে। একজন লোক $8\,\text{m\,s}^{-1}$ বেগে হাঁটলে ছাতা উল্লম্বের সাথে কত কোণে ধরতে হবে?}{সামনের দিকে উল্লম্বের সাথে $53.13^\circ$}
\topicbar{Topic-4: ডট ও ক্রস গুণন}
\problem[লেকচার ০১ · ডট গুণন]{$\vec A = 2\hat i + 3\hat j - \hat k$ এবং $\vec B = \hat i - 2\hat j + 2\hat k$ হলে এদের মধ্যবর্তী কোণ নির্ণয় করো।}{$\theta = 122.3^\circ$}
\problem[লেকচার ০১ · ক্রস গুণন]{উপরের $\vec A$ ও $\vec B$ এর জন্য $\vec A\times\vec B$ ও এর মান নির্ণয় করো।}{$4\hat i - 5\hat j - 7\hat k$; মান $\sqrt{90} = 9.49$ একক}
\problem[লেকচার ০১ · শর্ত]{$a$ এর মান কত হলে $2\hat i + a\hat j + \hat k$ এবং $4\hat i - 2\hat j - 2\hat k$ ভেক্টরদ্বয় পরস্পর লম্ব হবে?}{$a = 3$}
\problem[লেকচার ০১ · একক ভেক্টর]{$\vec P = 3\hat i + 4\hat j$ এর দিকে একক ভেক্টর নির্ণয় করো।}{$\hat p = \tfrac{1}{5}(3\hat i + 4\hat j)$}
\problem[লেকচার ০১ · ক্ষেত্রফল]{$\vec A = 3\hat i + \hat j + 2\hat k$ ও $\vec B = \hat i - \hat j + \hat k$ দ্বারা গঠিত সামান্তরিকের ক্ষেত্রফল কত?}{$\sqrt{26} = 5.10$ বর্গ একক}
\end{multicols*}
\end{document}
